% ============================================================
% MÓDULO 1: BANDIDOS MULTI-BRAZO
% ============================================================
\section{M1 · Bandidos}
\module{Módulo 1: Bandidos Multi-Brazo}{Aprender por evaluación en un escenario de decisión estático}

\begin{frame}{OptiBot y la Selección de Proveedores}
  \begin{columns}[c,onlytextwidth]
    \column{0.52\textwidth}
      \textbf{El Escenario en la Estación de Despacho:}
      \begin{itemize}
        \item OptiBot debe despachar paquetes eligiendo entre \textbf{4 proveedores} de mensajería: $\{A, B, C, D\}$.
        \item Cada proveedor tiene tiempos y costos variables (estocásticos).
        \item Cada envío genera una recompensa neta $R_t$.
        \item \textbf{Problema:} No conocemos cuál proveedor es el mejor de antemano.
      \end{itemize}
    \column{0.45\textwidth}
      \centering
      \begin{tikzpicture}[x=1.1cm,y=1.1cm]
        \node[draw=Ink,fill=Mist!30,rounded corners=4pt,inner sep=6pt,align=center] (robot) at (0,1.2)
          {\textbf{OptiBot}\\{\scriptsize Estación fija}};
        \node[draw=Teal,fill=Teal!12,rounded corners=3pt] (p1) at (2.4,2.2) {Prov. A ($q_*=3.2$)};
        \node[draw=Coral,fill=Coral!12,rounded corners=3pt] (p2) at (2.4,1.4) {Prov. B ($q_*=4.8$)};
        \node[draw=Gold,fill=Gold!15,rounded corners=3pt] (p3) at (2.4,0.6) {Prov. C ($q_*=2.1$)};
        \node[draw=Leaf,fill=Leaf!15,rounded corners=3pt] (p4) at (2.4,-0.2) {Prov. D ($q_*=4.1$)};
        \draw[-{Latex},thick,SoftText] (robot) -- (p1);
        \draw[-{Latex},thick,Coral,line width=1.5pt] (robot) -- (p2);
        \draw[-{Latex},thick,SoftText] (robot) -- (p3);
        \draw[-{Latex},thick,SoftText] (robot) -- (p4);
      \end{tikzpicture}
  \end{columns}
  \vspace{0.3cm}
  \keyidea{En Bandidos sólo hay \textbf{acciones} y \textbf{recompensas}. No hay estados: el entorno no cambia de lugar.}
\end{frame}

\begin{frame}{Valor de una Acción y el Dilema Central}
  \begin{columns}[T,onlytextwidth]
    \column{0.48\textwidth}
      \begin{block}{Valor Verdadero $q_*(a)$}
        \[
          q_*(a) \doteq \mathbb{E}[R_t \mid A_t = a]
        \]
        Es la ganancia esperada condicionada a elegir la acción $a$.
      \end{block}
      \vspace{0.2cm}
      Nuestra estimación en el paso $t$ se denota $Q_t(a)$. Queremos que $Q_t(a) \to q_*(a)$.
    \column{0.48\textwidth}
      \begin{alertblock}{Exploración vs. Explotación}
        \begin{itemize}
          \item \textbf{Explotar:} Elegir la acción con mayor $Q_t(a)$ actual (maximizar ganancia inmediata).
          \item \textbf{Explorar:} Probar acciones no óptimas para aprender su verdadero $q_*(a)$ (ganancia futura).
        \end{itemize}
      \end{alertblock}
  \end{columns}
  \vspace{0.4cm}
  \keyidea{El agente nunca puede hacer ambas cosas al mismo tiempo con una sola acción: debe alternar con criterio.}
\end{frame}

\begin{frame}{La Regla de Actualización Incremental}
  \begin{columns}[c,onlytextwidth]
    \column{0.52\textwidth}
      ¿Cómo actualizamos el promedio de una acción sin guardar todo su historial en memoria?
      \[
        Q_{n+1} = Q_n + \alpha \Big( \underbrace{R_n - Q_n}_{\text{Error de Predicción}} \Big)
      \]
      \begin{itemize}
        \item Si $\alpha = \frac{1}{n}$, es el promedio muestral exacto.
        \item Si $\alpha = \text{cte} \in (0, 1]$, pondera más los resultados recientes (entornos no estacionarios).
      \end{itemize}
    \column{0.45\textwidth}
      \begin{beamercolorbox}[rounded=true,sep=0.7em]{block body}
        \centering\bfseries
        Nuevo $\leftarrow$ Viejo $+$ Paso $\times$ [Objetivo $-$ Viejo]
      \end{beamercolorbox}
      \vspace{0.3cm}
      \small \textbf{¡El patrón universal de RL!}\\
      Esta estructura reaparece en:
      \begin{itemize}
        \item Diferencia Temporal TD(0)
        \item SARSA y Q-Learning
        \item Gradiente descendente en Redes Neuronales
      \end{itemize}
  \end{columns}
\end{frame}

\begin{frame}[fragile]{Estrategia $\varepsilon$-Greedy en Python}
  \begin{columns}[T,onlytextwidth]
    \column{0.62\textwidth}
\begin{lstlisting}[style=rlcode]
import numpy as np

k = 4  # 4 proveedores
q_true = np.array([3.2, 4.8, 2.1, 4.1])  # desconocidos
Q = np.zeros(k)  # estimaciones iniciales
N = np.zeros(k)  # conteo de elecciones
epsilon = 0.10

for t in range(1, 1001):
    # Exploracion vs Explotacion
    if np.random.rand() < epsilon:
        action = np.random.choice(k)     # Explorar
    else:
        action = np.argmax(Q)            # Explotar (greedy)
    
    # Simular recompensa con ruido normal
    reward = np.random.normal(loc=q_true[action], scale=1.0)
    
    # Actualizacion incremental
    N[action] += 1
    Q[action] += (1.0 / N[action]) * (reward - Q[action])

print("Q estimado:", np.round(Q, 2))
# Salida convergente: [3.21, 4.79, 2.12, 4.09]
\end{lstlisting}
    \column{0.35\textwidth}
      \begin{exampleblock}{La gran limitación de M1}
        OptiBot aprende a elegir al proveedor $B$ óptimamente.\\
        \vspace{0.2cm}
        \textbf{Pero...} sus decisiones no cambian el mundo. El almacén siempre está en el mismo estado.\\
        \vspace{0.2cm}
        ¿Qué pasa cuando las acciones cambian el estado y el entorno evoluciona?
      \end{exampleblock}
  \end{columns}
\end{frame}
